%% =============================================================================
%% LaTeX Template — JAMA Network Open Published-PDF Style
%% =============================================================================
%% Approximates the visual layout of JAMA Network Open research articles.
%% NOTE: JAMA Network Open has no official LaTeX class. This template
%% replicates the published PDF appearance for drafting / preprint use.
%% For submission, follow: https://jamanetwork.com/journals/jamanetworkopen/
%%                         pages/instructions-for-authors
%% =============================================================================

\documentclass[10pt, letterpaper]{article}

%% --- Packages ----------------------------------------------------------------
\usepackage[utf8]{inputenc}
\usepackage[T1]{fontenc}

% --- Fonts (approximating JAMA Network Open) ----------------------------------
% Body: serif (Times-like); Headings/sidebar: sans-serif (Helvetica-like)
\usepackage{mathptmx}          % Times serif for body
\usepackage[scaled=0.92]{helvet} % Helvetica for sans-serif elements
\usepackage{courier}           % Monospace fallback

\usepackage[
  top=1.1in, bottom=0.85in,
  left=0.75in, right=0.75in,
  headheight=30pt
]{geometry}

\usepackage{setspace}
\usepackage{graphicx}
\usepackage{booktabs}
\usepackage{colortbl}          % for \arrayrulecolor
\usepackage{array}
\usepackage{tabularx}
\usepackage{longtable}
\usepackage{multirow}
\usepackage{float}
\usepackage{caption}
\usepackage{subcaption}
\usepackage[numbers,super,comma,sort&compress]{natbib}
\usepackage{url}
\usepackage{xcolor}
\usepackage{enumitem}
\usepackage{titlesec}
\usepackage{fancyhdr}
\usepackage{lineno}
\usepackage{amsmath}
\usepackage{microtype}
\usepackage{ragged2e}
\usepackage{tikz}
\usepackage{wrapfig}
\usepackage{mdframed}
\usepackage{changepage}        % for adjustwidth
\usepackage{lastpage}
\usepackage[colorlinks=true,
  linkcolor=jamacrimson,
  citecolor=jamacrimson,
  urlcolor=jamacrimson]{hyperref}

%% === COLOR PALETTE (from published PDF) ======================================
\definecolor{jamacrimson}{RGB}{175, 30, 55}    % Deep red for "Original Investigation", links
\definecolor{jamagold}{RGB}{195, 155, 50}      % Accent gold (JAMA Open icon)
\definecolor{jamagray}{RGB}{90, 90, 90}        % Body gray
\definecolor{jamalightgray}{RGB}{200, 200, 200}% Rule lines
\definecolor{jamadarkgray}{RGB}{50, 50, 50}    % Headings, title

%% === PAGE STYLE ==============================================================
\setstretch{1.08}              % Tight leading like published PDF
\setlength{\parindent}{0pt}
\setlength{\parskip}{6pt}

%% --- Running header: crimson bar + journal/subject + paper title -------------
%% Store paper title and subject for headers
\newcommand{\jamashorttitle}{Running Title of the Article}
\newcommand{\jamasubject}{Public Health}

\pagestyle{fancy}
\fancyhf{}
\renewcommand{\headrulewidth}{0pt}
\renewcommand{\footrulewidth}{0.4pt}
% Running header on continuation pages: colored bar + journal | subject + title
\fancyhead[L]{%
  \smash{\makebox[0pt][l]{%
    \raisebox{10pt}{\textcolor{jamacrimson}{\rule{0.25\textwidth}{4pt}}%
      \textcolor{jamadarkgray}{\rule{0.75\textwidth}{4pt}}}}}%
  \vspace{6pt}\\
  {\sffamily\bfseries\fontsize{8}{10}\selectfont
    \textcolor{jamadarkgray}{JAMA Network Open\enspace|\enspace}%
    \textcolor{jamacrimson}{\jamasubject}}%
  \hfill
  {\sffamily\fontsize{7.5}{9.5}\selectfont\color{jamagray}\jamashorttitle}%
}
\fancyfoot[L]{%
  \footnotesize\textit{JAMA Network Open.}
  20XX;X(X):eXXXXXXX.
  doi:10.1001/jamanetworkopen.20XX.XXXXXXX}
\fancyfoot[R]{%
  \footnotesize Month Day, 20XX\qquad\thepage/\pageref*{LastPage}}

% First-page style (no header, special footer)
\fancypagestyle{firstpage}{%
  \fancyhf{}
  \renewcommand{\headrulewidth}{0pt}
  \renewcommand{\footrulewidth}{0.4pt}
  \fancyfoot[L]{%
    \footnotesize\textit{JAMA Network Open.}
    20XX;X(X):eXXXXXXX.
    doi:10.1001/jamanetworkopen.20XX.XXXXXXX
    \hfill (Reprinted)\hfill
    Month Day, 20XX\quad\thepage/\pageref*{LastPage}}
}

%% === SECTION FORMATTING ======================================================
% JAMA uses bold sans-serif headings in crimson
\titleformat{\section}
  {\sffamily\bfseries\fontsize{11}{13}\selectfont\color{jamacrimson}}
  {}{0em}{}
\titleformat{\subsection}
  {\sffamily\bfseries\fontsize{10}{12}\selectfont\color{jamacrimson}}
  {}{0em}{}
\titleformat{\subsubsection}
  {\sffamily\itshape\fontsize{10}{12}\selectfont\color{jamacrimson}}
  {}{0em}{}

\titlespacing*{\section}{0pt}{14pt}{4pt}
\titlespacing*{\subsection}{0pt}{10pt}{3pt}
\titlespacing*{\subsubsection}{0pt}{8pt}{2pt}

%% === CAPTION FORMATTING ======================================================
\captionsetup{
  font={small,sf},
  labelfont={bf},
  labelsep=period,
  justification=justified,
  singlelinecheck=false,
  skip=6pt,
  textfont={color=jamadarkgray}
}

%% === KEY POINTS BOX STYLE ====================================================
\mdfdefinestyle{keypointsbox}{%
  linewidth=0pt,
  innerleftmargin=10pt,
  innerrightmargin=10pt,
  innertopmargin=8pt,
  innerbottommargin=8pt,
  leftline=false, rightline=false,
  topline=true, bottomline=true,
  linecolor=jamalightgray,
  backgroundcolor=white
}

%% === ABSTRACT LABEL STYLE ====================================================
\newcommand{\abslabel}[1]{%
  {\vspace{4pt}\noindent\sffamily\fontseries{b}\selectfont\fontsize{8.5}{10}\selectfont #1\enspace}%
}
\newcommand{\absrule}{%
  \vspace{3pt}\noindent{\color{jamacrimson}\rule{0.62\textwidth}{0.4pt}}\vspace{3pt}%
}

%% === JAMA NETWORK OPEN LOGO PLACEHOLDER ======================================
\newcommand{\jamastyletopbar}{%
  \noindent
  \begin{minipage}[t]{0.55\textwidth}
    {\sffamily\bfseries\fontsize{16}{18}\selectfont
      \textcolor{jamadarkgray}{JAMA}\,|\,%
      \textcolor{jamadarkgray}{\fontseries{l}\selectfont Network}\,|\,%
      {\fontsize{20}{22}\selectfont\textcolor{jamadarkgray}{Open}}%
      \textcolor{jamagold}{\raisebox{2pt}{\tiny\texttrademark}}}
  \end{minipage}%
  \hfill
  \begin{minipage}[t]{0.1\textwidth}
    \raggedleft
    \textcolor{jamagold}{\scriptsize\sffamily\raisebox{0pt}{\fbox{\tiny OA}}}
  \end{minipage}
  \vspace{2pt}

  \noindent{\color{jamalightgray}\rule{\textwidth}{0.6pt}}
  \vspace{6pt}
}

%% =============================================================================
%%  DOCUMENT BEGINS
%% =============================================================================
\begin{document}
\thispagestyle{firstpage}
\sloppy


%% --- Article Type + Subject --------------------------------------------------
{\sffamily\bfseries\fontsize{9}{11}\selectfont
  \textcolor{jamacrimson}{Original Investigation}%
  \textcolor{jamadarkgray}{\enspace|\enspace Public Health}%
}

\vspace{8pt}

%% --- Title -------------------------------------------------------------------
{\sffamily\bfseries\fontsize{18}{21}\selectfont\color{jamadarkgray}
Title of the Article\par}

\vspace{8pt}

%% --- Authors -----------------------------------------------------------------
{\fontsize{9.5}{11.5}\selectfont\color{jamadarkgray}
Author Name; Author Name; Author Name; AI Author Name\par}
%% List the AI models like Claude, Gemini you used into the author list

\vspace{12pt}

%% =============================================================================
%%  ABSTRACT  +  KEY POINTS  (side-by-side layout)
%% =============================================================================

\noindent{\color{jamalightgray}\rule{\textwidth}{0.5pt}}
\vspace{6pt}

\noindent
\begin{minipage}[t]{0.62\textwidth}
%% --- Abstract (left column, ~62% width) --------------------------------------
{\sffamily\bfseries\fontsize{11}{13}\selectfont\color{jamacrimson} Abstract\par}
\vspace{4pt}

\abslabel{IMPORTANCE}
{\fontsize{9}{11.5}\selectfont
Provide the rationale for the study and why it adds to current knowledge.
This section should be 1--2~sentences.\par}

\absrule

\abslabel{OBJECTIVE}
{\fontsize{9}{11.5}\selectfont
State the precise objective or research question. Begin with an infinitive
phrase such as ``To determine whether\ldots'' or ``To compare\ldots''\par}

\absrule

\abslabel{DESIGN, SETTING, AND PARTICIPANTS}
{\fontsize{9}{11.5}\selectfont
Describe the basic study design, setting, and participant eligibility
criteria. State the dates of the study and the dates analyses were
conducted.\par}

\absrule

\abslabel{EXPOSURE}
{\fontsize{9}{11.5}\selectfont
Describe the primary exposure(s) or intervention(s).\par}

\absrule

\abslabel{MAIN OUTCOMES AND MEASURES}
{\fontsize{9}{11.5}\selectfont
State the primary outcome(s) and how they were measured. Describe the
analytic approach.\par}

\absrule

\abslabel{RESULTS}
{\fontsize{9}{11.5}\selectfont
Begin with key demographics of the study sample (N, mean age, sex).
Then summarise the main quantitative findings with effect sizes and
confidence intervals.\par}

\absrule

\abslabel{CONCLUSIONS AND RELEVANCE}
{\fontsize{9}{11.5}\selectfont
Provide a brief interpretation of the results and their implications
for clinical practice, policy, or future research.\par}

\vspace{4pt}
\noindent{\color{jamalightgray}\rule{0.62\textwidth}{0.4pt}}


\end{minipage}%
\hfill
%% --- Key Points (right column, ~33% width) -----------------------------------
\begin{minipage}[t]{0.33\textwidth}
\begin{mdframed}[style=keypointsbox]
{\sffamily\bfseries\fontsize{10}{12}\selectfont\color{jamacrimson}
Key Points\par}
\vspace{4pt}

{\sffamily\bfseries\fontsize{8.5}{10.5}\selectfont Question\enspace}%
{\fontsize{8.5}{10.5}\selectfont
State the research question clearly and concisely.\par}

\vspace{6pt}

{\sffamily\bfseries\fontsize{8.5}{10.5}\selectfont Findings\enspace}%
{\fontsize{8.5}{10.5}\selectfont
Summarise the main findings of the study.\par}

\vspace{6pt}

{\sffamily\bfseries\fontsize{8.5}{10.5}\selectfont Meaning\enspace}%
{\fontsize{8.5}{10.5}\selectfont
State the implications of the findings.\par}
\end{mdframed}



\end{minipage}

\newpage

%% =============================================================================
%%  INTRODUCTION
%% =============================================================================
\section*{Introduction}

Provide background and context for the study. Discuss the significance of the
problem, summarise relevant prior work, identify gaps in the literature, and
state the study objective. The introduction should be concise (typically 3--5
paragraphs) and end with a clear statement of the study's aim.

Example citation format: Prior work has examined this
topic.\textsuperscript{1,2}


%% =============================================================================
%%  METHODS
%% =============================================================================
\section*{Methods}

\subsection*{Data}
Describe the data source(s), study population, and sample selection criteria.
Mention institutional review board approval status. State that the study
followed the STROBE, CONSORT, or other appropriate reporting guideline.

\subsection*{Outcome Measures}
Define the primary and secondary outcome variables and how they were measured.

\subsection*{Statistical Analysis}
Describe the statistical methods in sufficient detail for a knowledgeable
reader to replicate the analysis. Specify the software used. State the
significance threshold and whether tests were 1- or 2-sided.


%% =============================================================================
%%  RESULTS
%% =============================================================================
\section*{Results}

Begin with a description of the study sample (Table).
Then present the main findings organised by research question or hypothesis.
Include effect estimates, confidence intervals, and \textit{P}~values.

\vspace{6pt}

%% --- Table Template (JAMA style: no vertical rules, thin horizontal rules) ------------
\begin{table}[H]
\centering
\sffamily\fontsize{8.5}{11}\selectfont
\captionsetup{font={small,sf,bf}, textfont={normalfont,sf,color=jamadarkgray}}
\caption{Descriptive Statistics of Sampled Health Care Workers}
\label{tab:descriptive}

\begin{tabularx}{\textwidth}{@{} >{\raggedright\arraybackslash}p{3.8cm}
  >{\centering\arraybackslash}X
  >{\centering\arraybackslash}X
  >{\centering\arraybackslash}X @{}}

\arrayrulecolor{jamadarkgray}
\toprule[0.8pt]

\textbf{Characteristic}
  & \textbf{All\newline(N~=~31\,142)}
  & \textbf{Mandate states\newline(n~=~12\,431)\textsuperscript{a}}
  & \textbf{Nonmandate states\newline(n~=~18\,711)\textsuperscript{b}} \\

\midrule[0.5pt]

\multicolumn{4}{@{}l}{\textbf{Sex}} \\
\quad Female           & 24\,294 (72.1) & 9554 (71.2) & 14\,740 (72.8) \\
\quad Male             & 6848 (27.9)    & 2877 (28.8) & 3971 (27.2) \\[4pt]

\multicolumn{4}{@{}l}{\textbf{Marital status}} \\
\quad Unmarried        & 11\,845 (39.3) & 5052 (41.7) & 6793 (37.4) \\
\quad Married          & 19\,297 (60.7) & 7379 (58.3) & 11\,918 (62.6) \\[4pt]

\multicolumn{4}{@{}l}{\textbf{Race}} \\
\quad Black            & 2927 (15.3) & 1210 (14.3) & 1717 (16.1) \\
\quad White            & 24\,951 (71.7) & 9401 (66.5) & 15\,550 (75.7) \\
\quad Other\textsuperscript{c} & 3264 (13.0) & 1820 (19.2) & 1444 (8.2) \\[4pt]

\multicolumn{4}{@{}l}{\textbf{Ethnicity}} \\
\quad Hispanic         & 2555 (12.8) & 1326 (15.8) & 1229 (10.5) \\
\quad Non-Hispanic     & 28\,587 (87.2) & 11\,105 (84.2) & 17\,482 (89.5) \\[4pt]

\multicolumn{4}{@{}l}{\textbf{Age, y}} \\
\quad 25--34           & 5700 (28.6) & 2300 (29.0) & 3400 (28.4) \\
\quad 35--49           & 13\,547 (39.3) & 5327 (38.5) & 8220 (39.9) \\
\quad 50--64           & 11\,895 (32.0) & 4804 (32.5) & 7091 (31.7) \\[4pt]

\multicolumn{4}{@{}l}{\textbf{Educational level}} \\
\quad Less than college degree & 10\,999 (51.3) & 4169 (49.5) & 6830 (52.6) \\
\quad College degree or higher & 20\,143 (48.7) & 8262 (50.5) & 11\,881 (47.4) \\[4pt]

\multicolumn{4}{@{}l}{\textbf{Annual household income, \$}} \\
\quad $<$35\,000       & 2831 (12.4) & 1020 (11.1) & 1811 (13.4) \\
\quad 35\,000--100\,000 & 9778 (32.0) & 3583 (29.4) & 6195 (34.0) \\
\quad $>$100\,000      & 12\,439 (32.3) & 5353 (35.4) & 7086 (29.9) \\
\quad Income unknown   & 6094 (23.3) & 2475 (24.0) & 3619 (22.7) \\[4pt]

\multicolumn{4}{@{}l}{\textbf{Vaccine uptake outcomes}} \\
\quad Ever vaccinated in premandate period\textsuperscript{d}
  & 12\,172 (84.2) & 5147 (87.5) & 7025 (81.6) \\
\quad Ever vaccinated in postmandate period\textsuperscript{e}
  & 16\,038 (89.9) & 6449 (93.4) & 9589 (87.2) \\
\quad Primary series completed in premandate period\textsuperscript{f}
  & 11\,978 (82.5) & 5067 (86.1) & 6911 (79.7) \\
\quad Primary series completed in postmandate period
  & 15\,944 (89.3) & 6411 (92.8) & 9533 (86.7) \\

\bottomrule[0.8pt]
\end{tabularx}

\vspace{4pt}
\begin{flushleft}
\fontsize{7.5}{9.5}\selectfont\rmfamily\color{jamagray}
Percentage estimates were weighted using survey weights.\\[1pt]
\textsuperscript{a}\,California, Colorado, Connecticut, Delaware, Illinois,
  Maine, Maryland, Massachusetts, New Jersey, New Mexico, New York, Oregon,
  Pennsylvania, Rhode Island, Washington, and Washington, DC.\\[1pt]
\textsuperscript{b}\,Alabama, Alaska, Arizona, Arkansas, Florida, Georgia,
  Indiana, Iowa, Kansas, Louisiana, Michigan, Missouri, Montana, Nebraska,
  New Hampshire, North Carolina, North Dakota, Ohio, Oklahoma, South Carolina,
  South Dakota, Tennessee, Texas, Utah, Vermont, Virginia, West Virginia,
  Wisconsin, and Wyoming.\\[1pt]
\textsuperscript{c}\,The Household Pulse Survey categorises race as Asian,
  Black, White, multiracial, and any other race. This study grouped Asian,
  multiracial, and any other race as ``other'' because of small sample sizes.\\[1pt]
\textsuperscript{d}\,Premandate period included survey weeks 31 to 33.\\[1pt]
\textsuperscript{e}\,Postmandate period included survey weeks 34 to 39.\\[1pt]
\textsuperscript{f}\,HCWs who completed or intended to complete the primary
  series of COVID-19 vaccination.
\end{flushleft}

\end{table}


%% =============================================================================
%%  DISCUSSION
%% =============================================================================
\section*{Discussion}

Interpret the main findings in context of prior literature. Discuss mechanisms,
clinical or policy implications, and how results compare with existing evidence.

\subsection*{Limitations}
Acknowledge study limitations candidly.


%% =============================================================================
%%  CONCLUSIONS
%% =============================================================================
\section*{Conclusions}

Provide a concise summary of the main findings and their relevance.


%% =============================================================================
%%  ARTICLE INFORMATION
%% =============================================================================
\vspace{6pt}
\noindent{\color{jamalightgray}\rule{\textwidth}{0.4pt}}
\vspace{4pt}





%% =============================================================================
%%  REFERENCES
%% =============================================================================

{\fontsize{8.5}{10.5}\selectfont
\bibliographystyle{vancouver}
\bibliography{references}
}


%% =============================================================================
%%  SUPPLEMENT (optional, starts on new page)
%% =============================================================================
\clearpage
\section*{Supplement 1}

\subsection*{eAppendix 1. Statistical Models and Methods Details}
Provide supplemental methodological details here, including the statistical models and methods used in the study.

\subsection*{eTable 1. Supplementary Table Title}
% Insert supplementary table here.

\subsection*{eFigure 1. Supplementary Figure Title}
% Insert supplementary figure here.


\end{document}